
\section{2. Related Work}

\subsection{2.1 Attention Variants}

\subsubsection{2.1.1 Multi-Head Attention (MHA)}

Introduced by Vaswani et al. (2017), MHA splits the hidden dimension into 
$H$ heads, each with its own $Q, K, V$ projections. This is the standard 
configuration for BERT, RoBERTa, and most encoder-only models.

\subsubsection{2.1.2 Multi-Query Attention (MQA)}

Shazeer (2019) proposed MQA, which shares a single $K, V$ pair across all 
query heads. This reduces the KV cache by a factor of $H$ but often 
degrades quality, especially for tasks requiring head diversity.

\subsubsection{2.1.3 Grouped-Query Attention (GQA)}

Ainslie et al. (2023) proposed GQA as a middle ground, evaluated on 
T5-based models. GQA is now standard in LLaMA-2, Mistral, Qwen, and 
Gemma. Our work re-implements GQA in a unified framework with extensive 
architecture support.

\subsection{2.2 KV Cache Optimization}

\subsubsection{2.2.1 Eviction}

H2O (Zhang et al., 2023) keeps attention sinks and heavy hitters. 
StreamingLLM (Xiao et al., 2024) preserves initial tokens. These methods 
are orthogonal to GQA.

\subsubsection{2.2.2 Quantization}

KIVI (Liu et al., 2024) quantizes KV to 2-bit. KVQuant (Hooper et al., 2024) 
uses per-channel scaling. GQA can be combined with quantization for 
multiplicative savings.

\subsubsection{2.2.3 Low-Rank Compression}

Low-rank KV cache compression (e.g., MLA in DeepSeek-V2) projects KV into 
a low-rank space. This is an orthogonal approach to reducing the number of 
heads.

\subsection{2.3 Parameter-Efficient Fine-Tuning}

LoRA (Hu et al., 2021) freezes base weights and trains low-rank adapters. 
We use LoRA as a recovery mechanism after GQA injection.

\subsection{2.4 Position of Our Work}

\begin{table}[H]
\centering
\small
\begin{tabular}{lll}
\toprule
\textbf{Axis} & \textbf{Original GQA (Ainslie et al.)} & \textbf{Our Implementation} \\
\midrule
Scope & T5-based enc-dec & Enc, dec, encoder-only, GQA-native \\
Weight transfer & Not detailed & 4 architectures \\
Fallback & Not specified & Returns None on failure \\
RoPE & Partial & Full support \\
Companion methods & None & Integrated with A3, AASVD, etc. \\
\bottomrule
\end{tabular}
\end{table}

\textbf{Key distinction.} Our work is a \textbf{production-hardened re-implementation}, 
not a new algorithmic contribution. We focus on robustness, architecture 
coverage, and integration with the broader compression toolkit.

\subsection{2.5 Comparison with Related Variants}
